% !TEX root = task-2-3.tex
% Editable task bank. Append new BankTask blocks at the end of the document.
\documentclass[12pt,a4paper]{article}
\usepackage{physics-textbook}
\usepackage{graphicx,needspace}
\usetikzlibrary{decorations.pathmorphing}
\newif\ifShowSolutions
\ifdefined\TutorEdition
  \ShowSolutionstrue
\else
  \ShowSolutionsfalse % Change to \ShowSolutionstrue to display solutions.
\fi
\newcounter{banktask}
\newsavebox{\BankTaskBox}
\newcommand{\BankTaskContent}[3]{%
  \noindent\textbf{\thebanktask.}\quad #2\par
  \ifShowSolutions
    \Needspace{50mm}% Keep the solution heading with its opening text or diagram.
    #3\par
    \ifdefined\TutorEdition
      \ifcsname TutorSource#1\endcsname
        \nopagebreak\smallskip{\small\csname TutorSource#1\endcsname\par}%
      \fi
    \fi
  \fi
}
\newcommand{\BankTask}[3]{%
  \refstepcounter{banktask}\par\medskip
  \sbox{\BankTaskBox}{\begin{minipage}{\linewidth}%
    \BankTaskContent{#1}{#2}{#3}\end{minipage}}%
  \ifdim\dimexpr\ht\BankTaskBox+\dp\BankTaskBox\relax>\dimexpr\textheight-20mm\relax
    \Needspace{35mm}\BankTaskContent{#1}{#2}{#3}%
  \else
    \Needspace{\dimexpr\ht\BankTaskBox+\dp\BankTaskBox\relax}%
    \noindent\usebox{\BankTaskBox}\par
  \fi
  \medskip
}
\begin{document}
\BookSubsection{2.3}{Основные формулы молекулярной физики}


% Source sdamgia-806: https://phys-ege.sdamgia.ru/problem?id=806
\BankTask{sdamgia-806}{
Из контейнера с твердым литием изъяли $4\,\text{моль}$ этого вещества. При этом число атомов лития в контейнере уменьшилось на

$1$) $4 \cdot 10 ^{23}$ 

$2$) $12 \cdot 10 ^{23}$ 

$3$) $24 \cdot 10 ^{23}$ 

$4$) $36 \cdot 10 ^{23}$
}{
\smallskip\noindent\textbf{Решение}\par
В одном моле любого вещества содержится $6 \cdot 10 ^{23}$ частиц. Таким образом, после изъятия $4\,\text{моль}$ вещества число атомов лития в контейнере уменьшилось на 

$4 \cdot 6 \cdot 10 ^{23} =24 \cdot 10 ^{23} .$

Ответ: $3$.
}

% Source sdamgia-3209: https://phys-ege.sdamgia.ru/problem?id=3209
\BankTask{sdamgia-3209}{
Какая масса у $3\,\text{моль}$ неона?

$1$) $12\,\text{г}$

$2$) $36\,\text{г}$

$3$) $60\,\text{г}$

$4$) $18 \cdot 10 ^{- 23} \text{г}$
}{
\smallskip\noindent\textbf{Решение}\par
Молярная масса неона равна $\mu _{\mathrm{Ne}}=20 \cdot 10 ^{- 3} \frac{\text{кг}}{\text{моль}} .$Следовательно $3\,\text{моль}$ неона имеют массу $3\text{моль}\!\!\!\! \cdot 20 \cdot 10 ^{- 3} \frac{\text{кг}}{\text{моль}} =60\text{г}.$

Ответ: $3$.
}

% Source sdamgia-3361: https://phys-ege.sdamgia.ru/problem?id=3361
\BankTask{sdamgia-3361}{
В баллоне находится $0{,}01\,\text{моль}$ газа. Сколько примерно молекул газа находится в баллоне?

$1$) $10 ^{21}$

$2$) $6 \cdot 10 ^{21}$

$3$) $10 ^{2}$

$4$) $6 \cdot 10 ^{24}$
}{
\smallskip\noindent\textbf{Решение}\par
В одном моль любого вещества находится приблизительно $N_A=6 \cdot 10 ^{23}$ молекул. Следовательно, в баллоне находится примерно $0{,}01 \cdot 6 \cdot 10 ^{23} =6 \cdot 10 ^{21}$ молекул.

Ответ: $2$.
}

% Source sdamgia-936: https://phys-ege.sdamgia.ru/problem?id=936
\BankTask{sdamgia-936}{
Плотность алюминия в $3$ раза больше плотности льда. В $1$ моле алюминия содержится:

$1$) атомов в $3$ раза больше, чем в одном моле льда;

$2$) столько же атомов, сколько в одном моле льда;

$3$) в $3$ раза меньше атомов, чем в одном моле льда;

$4$) на $12 \cdot 10 ^{23}$ атомов больше, чем в одном моле льда.
}{
\smallskip\noindent\textbf{Решение}\par
Количество вещества — это физическая величина, характеризующая количество однотипных структурных единиц, содержащихся в веществе. То есть эта величина «пересчитывает» количество молекул (для структурных веществ) и количество атомов (для «простых» веществ). 

В одном моле любого вещества содержится столько же структурных единиц вещества, сколько атомов содержится в $12$ граммах углерода. Следовательно, в одном моле алюминия содержится столько же атомов, сколько молекул в одном моле льда. Но молекула воды состоит из трех атомов, а значит, количество атомов в одном моле алюминия в $3$ раза меньше, чем количество атомов в одном моле льда.

Ответ: $3$.
}

% Source sdamgia-3769: https://phys-ege.sdamgia.ru/problem?id=3769
\BankTask{sdamgia-3769}{
В сосуде А находится $28\,\text{г}$ молекулярного азота, а в сосуде Б — $44\,\text{г}$ углекислого газа. В каком сосуде находится больше атомов?

$1$) в сосуде А

$2$) в сосуде Б

$3$) В сосудах А и Б содержится примерно одинаковое число атомов

$4$) в сосуде, объем которого больше
}{
\smallskip\noindent\textbf{Решение}\par
Это задание — удобный повод вспомнить, что атомы и молекулы — не одно и тоже. 

Определим, какое количество газов содержится в обоих сосудах. Поскольку молярная масса молекулярного азота равна $28\text{г}/\text{моль},$ а молекулярная масса углекислого газа — $44\text{г}/\text{моль},$ заключаем, что в обоих сосудах по одному моль вещества. Количество вещества — это физическая величина, характеризующая количество однотипных структурных единиц, содержащихся в веществе. То есть эта величина «пересчитывает» количество молекул. Если теперь вспомнить, что молекула азота $( N_2 )$ состоит из двух атомов, а молекула углекислого газа $( CO_2 )$ — из трех, то сразу ясно, что атомов в сосуде с углекислым газом приблизительно в полтора раза больше.

Ответ: $2$.
}

% Source sdamgia-30494: https://phys-ege.sdamgia.ru/problem?id=30494
\BankTask{sdamgia-30494}{
В сосуде объемом $V$ находится идеальный одноатомный газ массой $m$ при температуре $T$. ($\mu$ — молярная масса газа, $R$ — универсальная газовая постоянная, $N_{A}$ — число Авогадро) 

\TaskMatchingColumns{Физическая величина}{$A$)  концентрация газа\par Б)  давление газа}{Формула}{$1$)  $\frac{mRT}{\mu V}$\par $2$)  $\frac{\mu RT}{m V}$\par $3$)  $\frac{mN_A}{\mu V}$\par $4$)  $\frac{\mu N_A}{V}$}

\begin{center}\begin{tabular}{|c|c|}\hline
$A$ & Б\\\hline
\strut & \strut\\\hline\end{tabular}\end{center}
}{
\smallskip\noindent\textbf{Решение}\par
А) Концентрация по определению $n= \frac{N}{V} ,$ число молекул $N=\nu N_A= \frac{m}{\mu} N_A,$ тогда получаем $n= \frac{mN_A}{\mu V}$ (пункт $3$).

Б) Из уравнения Клапейрона — Менделеева $pV= \frac{m}{\mu} RT$ выражаем давление газа $p= \frac{mRT}{\mu V}$ (пункт $1$).

Ответ: $31$.
}

% Source sdamgia-7788: https://phys-ege.sdamgia.ru/problem?id=7788
\BankTask{sdamgia-7788}{
\TaskMatchingColumns{Физическая величина}{А)  Удельная теплота парообразования\par Б)  Молярная масса}{Размерность}{$1$)  $\frac{\text{м} ^{2}}{\text{с} ^{2} \cdot \text{К}}$\par $2$)  моль\par $3$)  $\frac{\text{м} ^{2}}{\text{с} ^{2}}$\par $4$)  кг/моль}

\begin{center}\begin{tabular}{|c|c|}\hline
$A$ & Б\\\hline
\strut & \strut\\\hline\end{tabular}\end{center}
}{
\smallskip\noindent\textbf{Решение}\par
По определению удельная теплота парообразования — это количество теплоты, которое нужно передать одному килограмму вещества, чтобы перевести его из жидкого состояния в газообразное при условии, что температура тела равна температуре кипения. Значит, размерность этой величины $\frac{\text{Дж}}{\text{кг}} = \frac{\text{Н} \cdot \text{м}}{\text{кг}} = \frac{\text{кг} \cdot \text{м}/c ^{2} \cdot \text{м}}{\text{кг}} = \frac{\text{м} ^{2}}{\text{с} ^{2}} .$

Молярная масса — масса одного моля вещества, измеряется в кг/моль.

Ответ: $34$.
}

% Source sdamgia-901: https://phys-ege.sdamgia.ru/problem?id=901
\BankTask{sdamgia-901}{
При температуре $250$ $K$ и давлении $1{,}5 \cdot 10 ^{5} \text{Па}$ плотность газа равна $2\text{кг}/\text{м} ^{3} \;.$ Какова молярная масса этого газа? (В килограммах на моль с точностью до десятитысячных.)
}{
\smallskip\noindent\textbf{Решение}\par
Воспользуемся уравнением Клапейрона — Менделеева: $pV=\nu RT.$ Вспоминая, что $\nu = \frac{m}{M} ,$ где $M$ — искомая молярная масса, а $m=\rho V,$ получаем: 

$M= \frac{\rho RT}{p} = \frac{2 \cdot 8{,}31 \cdot 250}{1{,}5 \cdot 10 ^{5}} \text{кг}/\text{моль}=27{,}7 \cdot 10 ^{- 3} \text{кг}/\text{моль}=0{,}0277\text{кг}/\text{моль}.$

Ответ: $0{,}0277\,\text{кг}$/моль.
}

% Source sdamgia-31410: https://phys-ege.sdamgia.ru/problem?id=31410
\BankTask{sdamgia-31410}{
В закрытом сосуде неизменного объема при комнатной температуре находилась смесь аргона и гелия (по $1$ молю каждого газа). Половину содержимого этого сосуда выпустили наружу, а затем добавили в сосуд еще $2$ моля гелия. Газы считать идеальными, температуру — постоянной.

$1$. Парциальное давление аргона уменьшилось.

$2$. Давление смеси газов в сосуде не изменилось.

$3$. Концентрация гелия увеличилась.

$4$. В начале опыта плотности газов были одинаковые.

$5$. В начале опыта внутренние энергии газов были одинаковые.
}{
\smallskip\noindent\textbf{Решение}\par
$1$. Верно. Количество вещества аргона уменьшилось в $2$ раза, следовательно, из уравнения Клапейрона — Менделеева при неизменных объеме и температуре давление аргона уменьшилось.

$2$. Неверно. Вначале количество вещества в сосуде было равно $2$ моля, а после — $3$ моля. Следовательно, давление смеси газов увеличилось.

$3$. Верно. Поскольку количество вещества гелия увеличилось, то число молекул $N=\nu N_A$ также увеличилось, следовательно, концентрация $n= \frac{N}{V}$ увеличилась.

$4$. Неверно. Плотность вещества равна $\rho= \frac{m}{V} = \frac{\nu M}{V} .$ Поскольку молярные массы газов отличаются, то и плотность газов будет различной.

$5$. Верно. Внутренняя энергия одноатомных газов $U= \frac{3}{2} \nu RT$ будет одинаковой.

Ответ: $135$.
}

% Source sdamgia-19683: https://phys-ege.sdamgia.ru/problem?id=19683
\BankTask{sdamgia-19683}{
В некотором процессе $1$-$2$, который проводился с одним молем идеального одноатомного газа, среднеквадратичная скорость молекул газа возросла от $u_{1}=350\,\text{м/с}$ до $u_{2}=380\,\text{м/с}$. При этом давление $p$ газа в процессе оставалось постоянным, и газ совершил работу $A=292\,\text{Дж}$. Чему равна молярная масса $\mu$ этого газа?
}{
\smallskip\noindent\textbf{Решение}\par
$1$. Согласно молекулярно-кинетической теории идеального газа, среднеквадратичная скорость его молекул связана с температурой газа формулой:

$u = \sqrt{\frac{3RT}{\mu}} ,$ откуда $T = \frac{\mu u ^{2}}{3R} .$

$2$. Работа газа при изобарическом (согласно условию) процессе равна $A = p ( V_2 - V_1 ) .$

$3$. Из уравнения Клапейрона — Менделеева для одного моля идеального газа получаем, что $p ( V_2 - V_1 ) = R ( T_2 - T_1 ) = A.$

$4$. Таким образом, из записанных уравнений следует, что $A = \frac{\mu}{3} \cdot ( u ^{2} _2 - u ^{2} _1 ) ,$ откуда:

$\mu = \frac{3A}{( u ^{2} _2 - u ^{2} _1 )} = \frac{3 \cdot 292}{30 \cdot 730} = 0{,}040\text{кг}/\text{моль} = 40\text{г}/\text{моль}.$

Ответ: $40\,\text{г}$/моль.
}

\end{document}
